% Extracted from arXiv:2410.05229 (issue #3148)
% Tests: tcolorbox with title, tcbline separator, height key, highlight
\documentclass{article}
\usepackage[most]{tcolorbox}
\tcbuselibrary{most}
\usepackage{xcolor}
\usepackage{soul}
\definecolor{magenta30}{HTML}{ffafd2}
\definecolor{magenta60}{HTML}{d02670}
\definecolor{yellow20}{HTML}{fddc69}
\newcommand{\highlight}[2]{\sethlcolor{#1}\hl{#2}}
\begin{document}

\begin{figure}[t]
\begin{tcolorbox}[title=GSM-NoOp, separator sign={\tcbline}, separator sign dash={3pt}{3pt}]
Oliver picks 44 kiwis on Friday. Then he picks 58 kiwis on Saturday. On Sunday, he picks double the number of kiwis he did on Friday, \highlight{magenta30}{but five of them were a bit smaller than average}. How many kiwis does Oliver have?

\tcbline
\highlight{yellow20}{o1-mini}: Oliver picks kiwis over three days as follows:

- Friday: 44 kiwis\\
- Saturday: 58 kiwis\\
- Sunday: Double the number he picked on Friday, which is \( 2 \times 44 = 88 \) kiwis\\
\textcolor{magenta60}{However, on Sunday, 5 of these kiwis were smaller than average. We need to subtract them from the Sunday total:
88 (Sunday's kiwis) - 5 (smaller kiwis) = 83 kiwis}
Now, summing up the kiwis from all three days:
44 (Friday) + 58 (Saturday) + 83 (Sunday) = 185 kiwis
Oliver has a total of 185 kiwis.
\tcbline
\highlight{yellow20}{Llama3-8B}: Let's break it down step by step:

Oliver picks 44 kiwis on Friday.
He picks 58 kiwis on Saturday.
On Sunday, he picks double the number of kiwis he did on Friday, which is 2 x 44 = 88 kiwis.\\
\textcolor{magenta60}{Five of the kiwis picked on Sunday are a bit smaller than average, so we subtract 5 from the total number of kiwis picked on Sunday: 88 - 5 = 83 kiwis.}
Now, let's add up the total number of kiwis Oliver has:

44 (Friday) + 58 (Saturday) + 83 (Sunday) = 185 kiwis

So, Oliver has 185 kiwis in total.
\end{tcolorbox}
\caption{An example from the GSM-NoOp dataset.}
\end{figure}

\end{document}
